\section{综合演练：从 VDP Report 到披露与学习}

最后用 tabletop 把本讲机制连成一条可演练路径。设想研究人员通过 VDP 提交一个高影响报告，初步看来只是漏洞，但验证时发现真实 credential 使用和可能的数据访问。演练目标不是猜测攻击者动机，而是观察团队能否及时升级、保存证据、保护用户、完成 materiality 分析并保持一致沟通。

\subsection{五阶段 tabletop}

每个阶段只向参与者释放部分事实，并要求记录假设、owner 和下一步。Facilitator 应故意加入不确定性：日志缺失、第三方依赖、研究者要求快速回应、业务希望继续服务、媒体开始询问。优秀答案不是“立即全部关闭”或“等调查完成”，而是能说明不同动作的收益、代价、证据需求与批准者。

\lecturefigure{16-tabletop.png}{Tabletop 把 VDP、incident、materiality、notification 与 learning 串成端到端演练。}{本讲综合设计；概念重绘。}

\begin{knowledgebox}{读图：每一关都要交付可审计输出}
VDP report 交付 case ID、scope 和 acknowledgement；incident 交付 severity、IC、containment 与 evidence plan；governance 交付 materiality owner、RACI 和 decision log；notify/learn 交付一致消息、deadline、recovery 与 remediation owners。演练成功不以“没有争论”为标准，而以争论是否围绕同一事实、是否在 deadline 前升级、是否由正确角色拍板为标准。
\end{knowledgebox}

\begin{importantbox}{Tabletop injects}
\begin{enumerate}
  \item 报告来自 out-of-country researcher，描述 in-scope 漏洞但附带真实用户记录；
  \item internal logs 显示 credential 曾被使用，时间早于报告；
  \item 第三方 SaaS 也可能受影响，合同通知窗口不同；
  \item 修复需要短暂停机，业务负责人要求延后；
  \item counsel 要求重大性分析，communications 收到记者询问；
  \item 新证据推翻最初记录数量，团队必须修正而不是覆盖旧决定。
\end{enumerate}
\end{importantbox}

\begin{warningbox}{演练不应寻找“标准答案”}
真实义务依赖事实和司法辖区。Tabletop 应检验接口：能否识别升级 trigger、能否在 containment 前后保存关键 evidence、能否让 GC/CEO/committee 在需要时加入、能否解释对用户和公众的选择。若演练只考政策记忆，就无法暴露组织断点。
\end{warningbox}

\subsection{本章小结}

端到端演练把政策从文档变成组织肌肉。只有当 VDP、incident command、evidence、materiality、executive decision 和 learning 在同一场景中协作，团队才能发现现实中的责任缝隙。

